\section{Local Model Release Packaging Standard}
\label{sec:local-model-release-standard}

\subsection{Product contract}

Tessaris is not a general model repository.  It provides a curated local-model
library: every listed model has a defined task role, a supported device profile,
an installed runtime, an integrity record, and measured quality and speed
evidence.  A customer must be able to insert an SD card, select a model in
Vault, choose \emph{Download to SD card}, and use the model in Boardroom without
manually arranging weights, choosing a quantisation, or configuring an
inference engine.

The customer package is deliberately different from the internal master
archive.  Customers receive only the payload required for their selected device
and release.  Tessaris retains the source checkpoint or exact input quantisation,
the reference/control representation, build recipes, test corpus and release
evidence needed to reproduce and audit that package later.

\subsection{Standard release unit}

Every new model, including a model released after this document, must be
represented by a versioned Tessaris release unit with the following contents.

\begin{enumerate}
  \item \textbf{Identity and provenance.}  Record the upstream publisher,
  upstream revision or download URL, licence, model identifier, byte count and
  SHA--256 for every retained input artifact.  Community conversions must name
  their converter and revision; an unofficial or ``uncensored'' derivative is
  never silently substituted for an official model.
  \item \textbf{Selected customer payload.}  Include the chosen weight
  representation, tokenizer, configuration, template and any required vision
  files.  Include AION-generated expert banks, shards or caches only when they
  are necessary for the selected runtime.  Do not include internal controls,
  abandoned variants, teacher captures or other experimental material in a
  customer download.
  \item \textbf{Runtime profile.}  Specify the engine and exact build,
  required native kernels, thread and memory limits, cache policy, minimum
  internal free space, SD-card reserve, supported architectures and minimum
  unified memory.  The profile must fail closed on an unsupported device rather
  than silently choose a weaker configuration.
  \item \textbf{Integrity card.}  Supply \texttt{model-package.v1.json} with
  hashes, relative paths, runtime profile, compatibility limits, licence notice
  and release status.  A production card must also carry a Tessaris release
  signature.  The present JSON card format and catalogue are implemented in
  \texttt{backend/modules/aion\_inference/model\_packages/}; development cards
  are not a substitute for a signed public release.
  \item \textbf{Acceptance evidence.}  Preserve a device-clean installation
  receipt, completed-answer quality suite, speed suite spanning ordinary
  business prompts, failure/fallback behaviour and reproducible command or
  runner version.  A best-case speed measurement is never the public speed
  claim.  The published claim is a range measured from the approved normal-use
  suite on the named hardware class.
\end{enumerate}

The installed card layout should be stable across all models:

\begin{verbatim}
TessarisModels/<model-id>/<release-version>/
  model-package.v1.json
  payload/                 # customer-selected weights and required assets
  runtime/                 # profile, compatible engine and native components
  receipts/                # install and acceptance receipts, not hidden claims
  LICENSES/                # upstream and Tessaris notices
\end{verbatim}

Vault reads the release catalogue before displaying a selectable model.  A
catalogue entry may be visible while it is being prepared, but it must remain
non-selectable until the package card, payload verification, hardware check and
acceptance result all pass.  Vault currently uses
\texttt{catalogue.v1.json} and \texttt{LocalModelVault} for this discovery
step.  The future ``Download to SD card'' operation must write the standard
layout atomically: download to a temporary folder, verify all hashes, write the
card and receipt, then make the release visible to Vault.

\subsection{Release procedure for every model}

\begin{enumerate}
  \item Acquire the official model and record provenance and hashes before
  optimisation.
  \item Produce one or more candidate representations.  Retain the exact
  source, control/reference representation and all deterministic build recipes.
  Optimisation changes must be separately identifiable; no candidate inherits a
  performance or quality claim merely because its source model passed one.
  \item Run correctness, task-quality and normal-use throughput acceptance on
  the intended hardware classes.  Preserve failed experiments as internal
  evidence, but do not ship them.
  \item Build the minimal customer payload and generate the package card,
  complete hash inventory, licence bundle and runtime profile.
  \item Perform a clean SD-card install using only that package.  Verify that
  Vault discovers it, Boardroom can start it, and removal of the card makes the
  model unavailable rather than simulated.
  \item Sign the release, publish the package to managed Tessaris storage, add
  the approved entry to the Vault catalogue, and publish only the speed and
  hardware range established by the clean acceptance receipt.
\end{enumerate}

This procedure makes future models routine: select a model for a defined role,
optimise it, certify it, package it in the same layout, publish it, and let
Vault install and recognise it.  Larger frontier models that do not pass the
local hardware and speed gates remain API-routed providers rather than being
forced into the SD catalogue.

\subsection{Current retained material and release state}

On 13 September 2026, the small-model master archive was created at
\begin{verbatim}
/Users/kevinrobinson/AION-Artifact-Archive/2026-09-13-small-model-master
\end{verbatim}
It is approximately 33~GiB.  It is an internal recovery archive, not a
customer package.  Its file \texttt{SHA256\_VERIFICATION.txt} records matching
source-card and archive SHA--256 values for the following material:
\begin{itemize}
  \item Qwen3--30B--A3B Q2 expert candidate:
  \texttt{7b1165eef9a868954a3fc9bdc31e3f454a04ebd86b14040945074c99edce5bac};
  \item Qwen3--30B--A3B Q4\_K\_M reference:
  \texttt{0d003f6662faee786ed5da3e31b29c978de5ae5d275c8794c606a7f3c01aa8f5};
  \item both Granite 3.1 3B A800M checkpoint tensors, together with its
  tokenizer and configuration files.
\end{itemize}

The original complete Granite BF16 expert-layer pack set is retained separately
at
\begin{verbatim}
/Users/kevinrobinson/AION-Artifact-Archive/
  2026-09-09-granite-layer-packs/granite-3.1-3b-a800m-v1
\end{verbatim}
The source code, native kernels, build recipes, technical records and Vault
integration live in
\begin{verbatim}
/Users/kevinrobinson/dev/COMDEX
\end{verbatim}
and the committed history was pushed to the \texttt{origin/main} Tessaris
repository on 13 September 2026.  Model binaries are intentionally not placed
in Git; they belong in the master archive and, once signed, managed release
storage.

The Qwen Q2 candidate has a recorded 50+ tokens/s result on its specialised
resident path, but it has not completed broad normal-business acceptance and is
therefore not a public speed claim or selectable Vault release.  Granite is
present and recoverable, but likewise requires a package card and full
acceptance.  GPT--OSS 120B remains a separate, acceptance-pending exact package;
its original SD card must be independently archived before that card is reused.

\subsection{Storage rule}

An SD card may be repurposed only after two independent conditions hold:
\begin{enumerate}
  \item every irreplaceable source and reference artifact has an independently
  stored, hash-verified master copy; and
  \item the release code, recipes, manifests, evidence and package definitions
  have been committed and pushed.
\end{enumerate}
After a signed package is placed in managed Tessaris storage, the internal Q4
control may be moved from the Mac to that storage if internal disk space is
needed.  It must remain recoverable for audit and rebuild; it is not discarded
merely because customers do not download it.

\subsection{Boardroom intelligence stack and role harnesses}

The SD-card library solves a storage and portability problem, not a working
memory problem.  It allows a customer to keep several approved local-model
packages without filling the computer's internal drive.  When AION activates a
model, the selected representation must still fit the device's available
unified memory or RAM/VRAM, including operating-system headroom and the chosen
context window.  Models that do not fit must not be loaded by swapping weights
through slow storage or by presenting a simulated response.

Consequently, an SD card is a portable library rather than an inherent source
of token-speed improvement.  An internal NVMe SSD is normally preferable for
initial loading on a desktop or Windows workstation.  Once the selected model
is resident in working memory, normal generation speed is primarily determined
by the model representation, memory bandwidth, execution runtime and hardware.
The SD package provides portability, easy installation, release integrity and
the ability to change the available local library without consuming internal
storage.

A complete Boardroom installation may combine the following layers:
\begin{enumerate}
  \item \textbf{Local LLM library.}  Two or three compatible packages stored
  on SD card or approved local storage: a light everyday worker, a stronger
  private general-work model, and, where useful, a reasoning or domain
  specialist.  Only the model selected for the current task needs working
  memory; the others remain stored in the library.
  \item \textbf{Frontier API tier.}  One or more connected providers such as
  OpenAI, Claude, Gemini, Kimi or DeepSeek for tasks beyond local hardware or
  local quality gates.  These models are not stored on the SD card.  Vault
  stores the connection configuration securely and AION records which provider
  served each request.
  \item \textbf{Image, video and voice capabilities.}  Approved local or API
  providers may supply visual creation, speech-to-text and text-to-speech.  The
  provider is selected by the same compatibility, privacy, quality, speed and
  cost rules as an LLM.
  \item \textbf{AION router.}  AION inspects the task and selects a compatible
  local model, a specialised harness, or an API provider.  It must account for
  the device profile, available memory, storage location, privacy requirement,
  expected latency, cost and required capability.  It may retain a model warm
  for related tasks or release it before activating another one.
\end{enumerate}

Each local package may also contain versioned \emph{role harnesses}.  A harness
is the reusable work layer around a model: role instructions, approved tools,
templates, task schemas, routing rules and evaluation checks.  It gives a
model a defined business function rather than exposing only a blank chat
window.  Examples include coding harnesses with repository, testing and review
workflows; business harnesses with document, spreadsheet and reporting
workflows; finance harnesses with structured extraction, reconciliation and
calculation checks; and sales, marketing or operations harnesses with their
respective templates and knowledge sources.

Harnesses are separately versioned from model weights.  Their package card
must state compatible model releases, required tools, data-access boundaries,
evaluation evidence and fallback behaviour.  This permits a business to keep a
small, legible intelligence stack while AION delegates individual tasks to the
appropriate local specialist or frontier provider.

\subsection{Device capability negotiation}

Tessaris must not assume that a customer uses an 18~GB Mac.  The 18~GB profile
is a reference acceptance target, not a universal release.  On first scan and
before every local-model installation, Vault records a non-sensitive device
capability snapshot: operating system, CPU architecture, available unified
memory or RAM, supported GPU runtime and VRAM, internal free storage, connected
SD or external storage, and supported acceleration backends.  No model is
offered merely because its files fit on an SD card.

Each model family may have several signed, separately tested variants.  For
example, a dense 27B model may provide a compact 18~GB-Mac profile, a higher
quality 36~GB-Mac profile, one or more Windows/NVIDIA GPU profiles, and an
API fallback.  Every variant has its own payload hashes, memory floor, context
limit, runtime flags, measured normal-use speed range and quality receipt.
They are different releases for compatibility and claims, even when they share
the same upstream model family.

Vault compares the device snapshot with these profiles and presents only the
safe choices.  It may recommend the best local tier, offer an SD-card download
to the selected storage location, or route the request to an approved API when
no local tier clears the memory, speed or quality gates.  AION receives the
selected profile as part of the model identity, allowing the router to choose
between local and frontier capability without requiring the business to
understand hardware, quantisation or GPU terminology.

\subsection{Initial curated model catalogue and provider plan}

The first Tessaris catalogue must be deliberately small.  It is a tested
library of useful roles, not a mirror of every model release.  A model enters
the catalogue only after its licence, source integrity, device variants,
normal-use speed, task quality and packaging receipt have passed the standard
release gate in this document.

\begin{center}
\begin{tabular}{p{0.18\linewidth}p{0.30\linewidth}p{0.42\linewidth}}
\textbf{Publisher} & \textbf{Candidate} & \textbf{Proposed Tessaris role} \\
\hline
OpenAI & GPT--OSS--20B & Private reasoning and tool use on compatible
16--18~GB-class devices.  GPT--OSS--120B is retained as a specialist-hardware
package, not a standard laptop product. \\
Meta & Muse Glimmer & Local agent, coding and workflow-automation candidate. \\
Google & Gemma family & Lightweight offline assistant for broadly compatible
hardware.  Gemma is the local open-model family; Gemini remains a frontier API
provider. \\
Mistral & Mistral Small~4 or Ministral~3 8B & Generalist, coding and
multilingual local candidate.  The 8B tier is especially appropriate for
ordinary laptops. \\
DeepSeek & DeepSeek--R1--Distill--Qwen--14B & Private reasoning candidate. \\
Qwen & Qwen3.8 27B & Local Pro business model candidate. \\
\end{tabular}
\end{center}

The initial supported Local Library is therefore organised by role rather than
by publisher:
\begin{enumerate}
  \item \textbf{Fast Local.}  Granite~3B, Gemma, or Ministral~3 8B for quick,
  routine offline work on a wide range of compatible devices.
  \item \textbf{Business Pro.}  Qwen3.8 27B, subject to a signed device
  variant and measured business-quality and normal-speed receipts.
  \item \textbf{Private Reasoning.}  GPT--OSS--20B or
  DeepSeek--R1--Distill--Qwen--14B, selected only when the device profile
  clears their memory and latency requirements.
  \item \textbf{Local Agent and Coding.}  Muse Glimmer, subject to the same
  compatibility and agent-safety acceptance process.
\end{enumerate}

Claude, Gemini, Grok, Kimi and the full DeepSeek frontier models are not
downloadable local packages in this initial library.  They belong in the
\textbf{Frontier API tier}.  AION may use an approved connected provider when
the task requires its capability, the customer allows external processing, and
the local device or local quality tier is unsuitable.  The receipt must name
the provider and model used.

The catalogue also needs media and voice capability, each represented as an
independently selected provider profile rather than as an assumed extension of
the text model:
\begin{enumerate}
  \item \textbf{Images.}  Maintain one or more approved image-generation and
  image-understanding providers.  A local image package may be offered only on
  devices with a validated accelerator and memory profile; otherwise AION uses
  an approved image API with a clear cost and privacy receipt.
  \item \textbf{Video.}  Begin with an approved video API, including MiniMax
  where its commercial terms and product integration are accepted.  Local video
  releases are workstation products and require their own GPU, licence, speed,
  reliability and output-quality acceptance evidence.  They must never be
  presented as ordinary laptop or SD-card packages merely because the weights
  can be stored externally.
  \item \textbf{Voice.}  Provide separate speech-to-text and text-to-speech
  choices.  A compact, validated local voice package can serve privacy-sensitive
  or offline use; a higher-quality or multilingual API profile can be selected
  when the customer permits it.
\end{enumerate}

This creates a complete business intelligence stack: a small set of local
packages for private and economical routine work, specialist role harnesses,
and connected frontier, image, video and voice services when they are the
better fit.  Vault presents the safe installed choices while AION performs the
routing decision from the task, policy, capability and device profile.
